\section{Integrated statement analysis}

\begin{qfmethod}[A credit-facility decision from integrated statement analysis]\label{meth:integrated-credit-facility-decision}
\qffield{Statement} The purpose is to decide whether, and on what terms, to offer Meridian a one-year revolver, not to assign a published agency rating. The lender has a USD 30 million single-name facility limit, requires minimum cash 50 and total interest-bearing debt/EBITDA no more than 2, and will not rely on the distorted reports. First correct recognition and classification, assess collateral and obligations, forecast operating cash and financing, stress the assumptions and compare feasible terms \citep[chap.~12, solutions~3--7, PDF p.~221; chap.~18, solutions~1--7, PDF pp.~255--256]{statement-analysis-workbook-2020}.
\qffield{Stress case} Relative to the 2027 base, sales are 880, COGS 60 percent, operating expense $132+0.06Sales=184.8$, depreciation 55 and opening-debt interest 17.6. Closing AR is 20 percent of sales (176), inventory 45 percent of COGS (237.6), AP 20 percent of COGS (105.6), accruals 25 percent of operating expense (46.2). At the hypothetical 25 percent current tax rate, NI is 70.95 and CFO 49.15. Committed capex 90, principal 20 and planned dividends 45 would leave cash $-18.35$ before financing, so preserving cash 50 needs a 68.35 draw. Even an unconditional 60 facility fails by 8.35 and exceeds the lender's limit.
\qffield{Alternatives} Declining credit entirely avoids exposure but leaves a viable operating business with a temporary funding need. A 30 facility without changing dividends is insufficient. Suspending dividends releases 45, reducing the required draw to 23.35. Year-end cash becomes 50, debt 223.35, equity $482.5+70.95=553.45$, and PP\&E 465; assets $50+176+237.6+465=928.6$ equal AP 105.6 plus accruals 46.2 plus debt 223.35 plus equity 553.45. EBITDA is 167.2, so leverage is 1.336. Further capex deferral could reduce need, but the specified 90 is committed and is not assumed away.
\qffield{Recommendation} Offer the 30 revolver conditional on dividend suspension in the stressed funding state, corrected reporting, verified receivable/inventory eligibility and monthly cash/covenant monitoring. It has 6.65 unused headroom in this stress and no base-case draw. Year-end draw timing avoids new interest this forecast year; carrying 23.35 into the next year adds 1.868 at 8 percent before repayment. Test intra-year seasonal cash needs before execution: an annual closing balance alone cannot prove daily liquidity. Decline or reduce the offer if eligibility or seasonal cash evidence invalidates the assumptions. The recommendation meets stated constraints under specified scenarios, not every possible downturn.
\end{qfmethod}

\begin{qfmethod}[Choosing and reconciling comparability adjustments]\label{meth:statement-comparability-adjustments}
\qffield{Statement} Compare firms on consistent reporting periods, consolidation, measurement and operating definitions. Identify the policy or estimate difference, decide whether a reliable adjustment is possible, quantify its income/balance/cash/tax effects and explain how the adjusted measure answers the decision. Do not force comparability with missing data or eliminate real economic differences such as a different product mix \citep[chap.~12, solutions~7--13, PDF pp.~221--222]{statement-analysis-workbook-2020}.
\qffield{Practice} Peer Northstar reports LIFO inventory 500, a closing LIFO reserve 200 and opening reserve 120. For a FIFO analytical comparison, add 200 to inventory and subtract the 80 reserve increase from COGS. Assuming unchanged tax payments and a 25 percent tax rate, add deferred-tax liability 50 and equity 150 at the reporting date; current-period profit rises 60 after tax and deferred tax expense 20. CFO is unchanged: higher income is offset by the extra working-capital accumulation and deferred tax. Do not substitute the entire reserve for its period change in the income statement.

For a development-capitalizing IFRS peer, compare to a US expensing basis only when project costs, accumulated amortization and tax bases are available: remove the remaining development asset, reverse its related DTL and adjust opening equity and current expense consistently. For leases, separate expense-pattern differences from liabilities already recognized; adding the same liability twice overstates leverage. For pensions, compare service cost and funded risk rather than taking reported periodic expense at face value.
\qffield{Recommendation} Use both reported and adjusted ratios, with a bridge and unresolved ranges. In Meridian's case the sale and maintenance corrections are supported by stipulated transaction facts and are required before lending; an estimated intangible value from management without support should not be inserted to inflate collateral. A consistent conservative lending haircut can coexist with an equity valuation that recognizes an intangible's earning power; the decision purpose differs.
\end{qfmethod}

\begin{qfmethod}[Accounting changes, normalization and balance-sheet modifications]\label{meth:accounting-change-statement-bridges}
\qffield{Statement} Distinguish an accounting-standard transition, a voluntary policy change, a new estimate, an error correction and an analyst's normalization. Each has different comparative-period treatment. Specify the journal-style balance-sheet effects first, then income/OCI and cash consequences, and only then recompute ratios. A reclassification changes presentation; it does not necessarily change total income, equity or cash \citep[IAS~8 paras.~19, 22, 34, 36 and~42; IAS~16 para.~39; IAS~12 para.~61A]{adopted-accounting-standards-2023}.
\qffield{Practice} An IFRS PP\&E revaluation increases an asset 20; at a hypothetical 25 percent tax rate, DTL rises 5 and OCI/equity 15, with no cash or current P\&L change. Higher future depreciation can lower future earnings. By contrast, the useful-life revision in the expense entry changes future depreciation without increasing the asset or restating prior years. Meridian's maintenance correction removes asset 20 and DTL 5 and reduces equity 15, while decreasing CFO by 20 and increasing investing cash flow by 20: the payment moves from investing to operating, with total cash unchanged; the reporting-quality table shows the full combined bridge. Adjusting a transitory disposal gain out of normalized income leaves actual disposal cash in investing flows.
\qffield{Transition comparison} IFRS 18's retrospective presentation transition can change subtotals and cash-flow classifications without improving the business. US ASU 2015-03 removed debt-issuance-cost assets and reduced the reported debt carrying amount by the same sum on retrospective presentation: equity and contractual debt service are unchanged, but asset-based ratios change. US ASU 2025-06 offers different internal-use software transition methods from its applicable period; modified transition can remove in-process capitalized costs that fail the new threshold. Do not describe every standards adoption as prospective or apply a future standard automatically to 2026 \citep[IFRS~18 paras.~C1--C3]{eu-2026-338-ifrs18}; \citep[Summary pp.~1--2]{us-debt-costs-2015}; \citep[Summary pp.~3--4]{us-software-2025}.
\qffield{Interpretation} After an impairment or buyback, apparent ROA/ROE improvement may reflect a smaller denominator. Retain opening-equity corrections separately from current profit and show unchanged total cash when only classification changes. Reassess the credit conclusion with corrected collateral and cash generation, rather than treating a mechanically improved ratio as economic recovery.
\end{qfmethod}
